\documentclass[11pt]{article}

\usepackage[margin=1in]{geometry}
\usepackage{amsmath,amssymb,amsfonts,amsthm,bm}
\usepackage{mathtools}
\usepackage{booktabs}
\usepackage{enumitem}
\usepackage[backend=biber,style=numeric,sorting=none,natbib=true]{biblatex}
\usepackage{hyperref}
\usepackage{xcolor}
\usepackage{multirow}

\hypersetup{
    colorlinks=true,
    linkcolor=blue,
    citecolor=blue,
    urlcolor=blue
}

\addbibresource{references.bib}

\title{Towards Certifiable Data Collection for Robotic Calibration}
\author{Parth Mahajan}
\date{\today}

\begin{document}

\maketitle

% -------------------------------------------------------
\begin{abstract}
Calibration is a foundational requirement for robotic perception systems,
yet calibration quality depends not only on the estimator but also on
the quality of the collected data.
In practice, calibration datasets are often gathered using heuristics such
as field-of-view coverage, motion diversity, or practitioner intuition,
while most existing pipelines focus on parameter estimation after data
collection.
This leaves open an important question: how can one certify that a chosen
set of calibration views, trajectories, or motion segments is sufficiently
informative for accurate parameter recovery?
In this work, we formulate calibration data collection as a certifiable
experimental design problem.
We cast the selection of calibration measurements as a subset selection
problem under an information-theoretic objective that strengthens the weakest
observable directions of the calibration parameter space.
Although the resulting optimization problem is NP-hard in general, we show
how greedy selection together with convex relaxation based post-hoc
verification can enable the efficient recovery of certifiably informative
calibration datasets in practice.
Unlike standard calibration pipelines such as Kalibr, which assume a
collected dataset and focus on estimation, and unlike next-best-view methods
that primarily optimize coverage or local uncertainty reduction, our
formulation targets the informativeness of the collected data itself.
We also position the method relative to recent learning-based calibration
work, which improves automation and targetless operation but generally does
not certify observability or data quality.
The resulting viewpoint suggests that calibration should be treated not only
as an estimation problem, but also as a certifiable data design problem.
\end{abstract}

% -------------------------------------------------------
\section{Introduction}
Calibration is a foundational requirement for modern robotic perception
systems.
Accurate estimates of camera intrinsics, sensor extrinsics, temporal
offsets, and other latent parameters are essential for downstream tasks such
as localization, mapping, sensor fusion, and control.
Although substantial progress has been made in calibration solvers and
toolchains, the success of these methods still depends strongly on the
quality of the collected calibration data.
Poorly chosen views, trajectories, or motion patterns can leave important
parameter directions weakly observable, resulting in unstable estimates and
degraded downstream performance.

In current practice, calibration data collection is often guided by simple
heuristics such as target coverage, motion diversity, or practitioner
intuition.
These strategies can be effective, but they do not explicitly reason about
whether the resulting dataset is sufficiently informative for accurate
parameter recovery.
Classical calibration pipelines such as Kalibr are highly effective at
estimating multi-camera and visual-inertial calibration parameters from a
provided dataset, but they largely assume that the dataset has already been
collected and do not directly solve the upstream problem of selecting
certifiably informative measurements \citep{furgale2013kalibr,kalibr_toolbox}.
Likewise, foundational camera calibration methods such as Zhang's
planar-pattern formulation provide highly practical and influential estimation
procedures, but not a principled notion of optimal data selection
\citep{zhang2000flexible}.

A related line of work comes from next-best-view and active calibration
methods, which select sensing actions or poses to improve visibility, reduce
uncertainty, or maximize information gain.
Systems such as AprilCal and Calibration Wizard explicitly acknowledge that
calibration quality depends strongly on the acquired data and therefore guide
the user toward improved views \citep{richardson2013aprilcal,peng2019calibration}.
Later work on efficient pose selection, next-best-view planning for
eye-in-hand calibration, and observability-aware active trajectory optimization
further reinforces the importance of calibration-aware data acquisition
\citep{rojtberg2018efficient,yang2023nbv,wang2025active}.
However, these methods are typically framed as sequential next-best-pose or
trajectory-planning procedures rather than as certifiable global
subset-selection problems over candidate calibration data.

Another growing research direction is deep learning-based calibration,
especially for targetless, online, or weakly supervised settings.
Methods such as RegNet, CalibNet, and LCCNet reduce manual effort and can
directly predict extrinsic corrections from multimodal data
\citep{schneider2017regnet,iyer2018calibnet,lv2021lccnet}.
More recent work continues to improve automation and online refinement
\citep{liao2023survey,huang2025whatmatters}.
At the same time, recent survey and critique papers note persistent challenges
in robustness, generalization, and consistent evaluation across sensor types,
scenes, and camera models \citep{liao2023survey,huang2025whatmatters}.
As a result, classical geometric pipelines remain the dominant choice in
precision-oriented robotic systems, even as learning-based methods continue
to broaden the scope of what can be automated.

In this paper, we cast calibration data collection as a certifiable
experimental design problem.
Rather than treating calibration purely as a downstream estimation task, we
formulate the selection of calibration views, poses, or motion segments as a
subset selection problem over candidate measurements.
Our objective is to strengthen the weakest observable directions of the
calibration parameter space through an E-optimal design criterion, and to
complement greedy selection with convex relaxation-based post-hoc verification.
This yields a principled framework for collecting calibration data that is
not only informative in practice, but certifiably informative under the
adopted information model.

The main contributions of this framing are:
\begin{enumerate}[leftmargin=1.5em]
    \item We position robotic calibration as a data design problem, not only
    as a downstream estimation problem.
    \item We connect calibration data selection to information-theoretic
    subset selection with emphasis on the weakest observable directions
    through E-optimality.
    \item We bridge active calibration with certifiable design by importing
    the greedy-plus-verification viewpoint from recent certifiable selection
    literature.
\end{enumerate}

% -------------------------------------------------------
\section{Related Work}

\subsection{Classical Calibration Backends}
Classical calibration methods estimate sensor intrinsics, extrinsics, and
timing parameters from a collected dataset.
Zhang's planar-pattern method remains a foundational reference for practical
camera calibration because it combines a closed-form initialization with
nonlinear refinement under a maximum-likelihood objective
\citep{zhang2000flexible}.
In robotics, Kalibr has become one of the most widely used practical backends
for multi-camera and visual-inertial calibration, supporting spatial and
temporal calibration together with rich camera model support
\citep{furgale2013kalibr,kalibr_toolbox}.
These methods are highly effective at estimation, but they primarily operate
on data that has already been collected.
Consequently, they do not directly answer the upstream question of how to
collect a calibration dataset that is certifiably informative.

\subsection{Self-Calibration, Observability, and Informative Data Selection}
A second line of work studies calibration through the lens of observability
and information content.
Maye et al.\ introduced a generic self-supervised calibration framework that
used information-theoretic measures to identify and retain novel measurement
sequences while handling unobservable directions explicitly
\citep{maye2013self}.
Schneider et al.\ later showed that visual-inertial self-calibration can be
posed as a segment-based problem over a small subset of informative motion
segments, yielding accuracy comparable to full-batch approaches at much lower
computational cost \citep{schneider2017informative}.
Their subsequent observability-aware work further emphasized that calibration
quality depends strongly on motion excitation and proposed information metrics
to assess which trajectory segments are worth retaining for calibration
\citep{schneider2019observability}.
Related work in online visual-inertial self-calibration provides a more
complete observability analysis and identifies degenerate motion primitives
that render subsets of calibration parameters unobservable
\citep{yang2023online}.
This literature strongly supports the premise that calibration failure is
often caused by weak observability and poor excitation, not merely by the
optimizer.

\subsection{Guided Capture, Next-Best-View, and Active Calibration}
Another closely related direction explicitly guides the collection process.
AprilCal introduced an assisted and repeatable camera calibration workflow
aimed at reducing reliance on operator intuition by suggesting improved
target placements and providing intuitive uncertainty summaries
\citep{richardson2013aprilcal}.
Calibration Wizard pushed this idea further by computing, from previously
acquired images, the next pose expected to yield the largest reduction in
intrinsic uncertainty while also accounting for corner-detection uncertainty
\citep{peng2019calibration}.
Additional work on efficient pose selection, next-best-view selection for
robot eye-in-hand calibration, and observability-aware active trajectory
optimization continues this trend
\citep{rojtberg2018efficient,yang2023nbv,wang2025active}.
These methods show that active calibration can significantly reduce the number
of required views or improve parameter uncertainty.
However, they are usually formulated as sequential guidance or local planning
procedures and do not provide a certificate that the final selected
calibration set is globally optimal or near-optimal under a global
subset-selection objective.

\subsection{Information-Theoretic Online Calibration}
Information-theoretic calibration ideas have also been explored in online and
multi-camera settings.
Dexheimer et al.\ developed an information-theoretic framework for online
multi-camera extrinsic calibration that used entropy-based keyframe selection
and segment evaluation to keep the estimation problem tractable
\citep{dexheimer2022information}.
These methods are conceptually close to the present viewpoint because they
explicitly distinguish informative from redundant data.
Nevertheless, they typically use heuristics such as entropy or keyframe
scoring rather than a certifiable E-optimal subset-selection formulation over
candidate calibration measurements.

\subsection{Certifiable Calibration Solvers}
There is also a growing literature on certifiable calibration solvers.
Giamou et al.\ formulated extrinsic calibration from per-sensor egomotion as
a quadratically constrained quadratic program and solved it via semidefinite
relaxation to obtain certifiably globally optimal estimates
\citep{giamou2019certifiable}.
Follow-on work extended this line to online-capable variants and related
hand-eye settings.
This literature is highly relevant conceptually, but its certification target
is different from ours: these methods certify the solution of a fixed
calibration estimation problem, whereas our goal is to certify the
informativeness of the selected dataset itself.
In other words, certifiable solver work addresses \emph{how to solve} the
posed calibration problem globally, while our framing addresses \emph{how to
collect} data that makes calibration well-conditioned in the first place.

\subsection{Targetless Geometric Calibration}
Targetless geometric calibration provides another important neighboring
direction, especially for cross-sensor calibration.
A classical example is the mutual-information-based LiDAR--camera calibration
method of Pandey et al., which aligns sensor observations without fiducial
targets and even reports a Cram\'er--Rao lower bound analysis
\citep{pandey2015mutual}.
This line of work demonstrates that calibration can be carried out directly
from natural scenes, making in-situ recalibration practical.
Yet these methods still optimize a registration objective over the data at
hand; they typically do not ask whether the selected data is globally optimal
under a calibration-specific experimental design objective.

\subsection{Learning-Based Calibration}
A final major direction is learning-based calibration.
Early milestones such as RegNet cast multimodal registration as direct deep
regression of the six-degree-of-freedom extrinsic transform
\citep{schneider2017regnet}, while CalibNet introduced self-supervised
geometric training without direct supervision on the calibration parameters
\citep{iyer2018calibnet}.
LCCNet later improved online LiDAR--camera self-calibration by exploiting
cost-volume style feature matching and iterative refinement \citep{lv2021lccnet}.
More recent surveys show that learning-based calibration has expanded
substantially across standard camera calibration, distortion estimation,
cross-view calibration, and cross-sensor calibration \citep{liao2023survey}.
At the same time, recent critique work argues that many regression-style
LiDAR--camera calibration networks still struggle in complex real scenes
\citep{huang2025whatmatters}.
For this reason, learning-based calibration is best viewed as complementary
to, rather than a replacement for, high-accuracy geometric calibration.
Most importantly for our setting, these methods improve automation and
targetless operation, but they generally do not certify observability or
the informativeness of the collected data.

\subsection{Relation to OASIS}
The work most closely aligned with our certifiable-design perspective is
OASIS, which formulates sensor arrangement for SLAM as a subset selection
problem under an E-optimality criterion and combines greedy selection with
convex-relaxation-based post-hoc verification \citep{kaveti2023oasis}.
OASIS shows that robotics design problems can be treated as
information-theoretic subset selection problems with meaningful certificates.
Our work borrows that conceptual template, but shifts the focus from sensor
arrangement for SLAM to the selection of informative views, poses, or motion
segments for calibration.
Thus, unlike Kalibr-style pipelines that estimate from a fixed dataset,
unlike AprilCal- and Calibration-Wizard-style methods that guide data
acquisition sequentially, and unlike certifiable solver approaches that
certify only the downstream optimizer, our goal is to certify the
informativeness of calibration data itself.

% -------------------------------------------------------
\section{Problem Statement}

In this section, we formulate calibration data collection as a subset
selection problem over a finite set of candidate measurements.
Our goal is to select a fixed-budget subset of calibration views, poses, or
motion segments that yields the most informative estimation problem for the
calibration parameters of interest.

\subsection{Parameterizing the Calibration Design Space}

Let
\begin{equation}
\mathcal{C} = \{c_1, c_2, \dots, c_N\}
\end{equation}
denote a finite set of $N$ candidate calibration data units, where each
candidate $c_i$ may represent a camera view, a robot pose, a trajectory
segment, or a batch of measurements.
We represent a selection by a binary vector $u \in \{0,1\}^N$.

\subsection{Calibration Likelihood Under a Selected Dataset}

Let $\theta \in \mathbb{R}^p$ denote the calibration parameters of interest
and $\eta$ any nuisance variables (target poses, robot poses, etc.).
For a selected subset $u$, the joint likelihood is
\begin{equation}
p(Z \mid \theta,\eta;u)
= \prod_{i=1}^{N} p_i(z_i \mid \theta,\eta)^{u_i}.
\end{equation}
The maximum-likelihood calibration estimate is
\begin{equation}
(\hat{\theta}(u),\hat{\eta}(u))
= \arg\min_{\theta,\eta} \sum_{i=1}^{N} u_i \, \ell_i(\theta,\eta),
\quad
\ell_i := -\log p_i(z_i \mid \theta,\eta).
\end{equation}

\subsection{Fisher Information and Marginal Calibration Information}

The Fisher information of candidate $c_i$ is
\begin{equation}
\mathcal I_i(\theta,\eta)
:= \mathbb E_{z_i}\!\left[-\nabla^2_{\theta,\eta}\log p_i\right],
\end{equation}
and the total FIM for a selection $u$ is additive:
$\mathcal I(\theta,\eta;u) = \sum_{i} u_i \mathcal I_i(\theta,\eta)$.

Partitioning into calibration parameters $\theta$ and nuisance $\eta$:
\begin{equation}
\mathcal I(u) =
\begin{bmatrix}
\mathcal I_{\theta\theta}(u) & \mathcal I_{\theta\eta}(u) \\
\mathcal I_{\eta\theta}(u) & \mathcal I_{\eta\eta}(u)
\end{bmatrix},
\end{equation}
the marginal calibration information matrix is the Schur complement
\begin{equation}
\widetilde{\mathcal I}_{\theta}(u)
= \mathcal I_{\theta\theta}(u)
- \mathcal I_{\theta\eta}(u)\,\mathcal I_{\eta\eta}(u)^{-1}\,\mathcal I_{\eta\theta}(u).
\label{eq:schur}
\end{equation}
Because the nuisance block $\mathcal I_{\eta\eta}$ is block-diagonal (each
candidate contributes an independent $6\times6$ pose block),
$\mathcal I_{\eta\eta}^{-1}$ costs $O(N \cdot 6^3)$---a cheap per-block
inversion---yielding a compact $p\times p$ matrix regardless of $N$.

\subsection{E-Optimal Calibration Data Selection}

The Cram\'er--Rao bound implies that the covariance of any unbiased estimator
is at least $\widetilde{\mathcal I}_{\theta}^{-1}$.
To strengthen the weakest observable direction we adopt the E-optimality
criterion:
\begin{equation}
\max_{u \in \{0,1\}^N,\;\mathbf{1}^\top u = K}
\;\lambda_{\min}\!\bigl(\widetilde{\mathcal I}_{\theta}(u)\bigr).
\label{eq:e_opt}
\end{equation}
This formulation turns calibration into an information-driven design problem:
rather than assuming all measurements are equally useful, we seek the subset
that provides the strongest worst-case observability.

% -------------------------------------------------------
\section{Experimental Design}
\label{sec:experiments}

We evaluate the proposed calibration data selection framework along four axes:
(i)~whether selected subsets improve calibration accuracy relative to
standard baselines,
(ii)~whether the E-optimal objective correlates with empirical estimation
quality,
(iii)~whether the method remains effective under weakly observable or
degenerate data regimes, and
(iv)~whether the post-hoc verification procedure provides meaningful
certificates on subset quality.
Since our goal is to study \emph{data design} rather than propose a new
calibration solver, all competing subset-selection methods are evaluated
using the same downstream calibration backend.

\subsection{Baselines}

We compare the following five methods:
\begin{enumerate}[leftmargin=1.5em]
    \item \textbf{Random:} uniform random $K$-subset (25 independent trials;
    mean reported).
    \item \textbf{Coverage:} greedy farthest-point sampling on camera
    translations in 3-D space.
    \item \textbf{Motion-diversity:} greedy farthest-point sampling on SO(3)
    rotations via geodesic distance.
    \item \textbf{Greedy-E:} greedy maximization of
    $\lambda_{\min}(\widetilde{\mathcal I}_{\theta})$.
    \item \textbf{FW-E (proposed):} Frank-Wolfe continuous relaxation of
    \eqref{eq:e_opt}, rounded to a binary solution; the FW duality gap
    provides a post-hoc certificate.
\end{enumerate}

\subsection{Metrics}

\paragraph{Information quality.}
$\lambda_{\min}(\widetilde{\mathcal I}_{\theta})$ (E-optimality, $\uparrow$).

\paragraph{Calibration quality.}
Because a single $\ell_2$ norm mixes parameters with different units and
physical interpretations, we decompose calibration error into three
physically interpretable groups:
\begin{itemize}[leftmargin=1.5em]
  \item $\epsilon_f$ \textbf{(focal-length error, px)}: $\sqrt{(\hat{f}_x - f_x^*)^2 + (\hat{f}_y - f_y^*)^2}$, measuring how well the image scale is recovered.
  \item $\epsilon_{pp}$ \textbf{(principal-point error, px)}: $\sqrt{(\hat{c}_x - c_x^*)^2 + (\hat{c}_y - c_y^*)^2}$, measuring optical-axis centering.
  \item $\epsilon_d$ \textbf{(distortion error, unitless)}: $\|\hat{\mathbf{d}} - \mathbf{d}^*\|_2$ where $\mathbf{d} = [k_1, k_2, p_1, p_2, k_3]^\top$, measuring lens-model accuracy.
\end{itemize}
\textbf{Train-RMS}: reprojection RMS on selected $K$ views (px).
\textbf{Held-out RMS}: mean per-point reprojection error on the remaining
$N-K$ views evaluated with estimated intrinsics and ground-truth poses (px).

\subsection{Setup}

All experiments use a pool of $N = 720$ candidate camera poses generated in
NVIDIA Isaac Sim around a fixed calibration target, with budget $K = 20$.
Both targets are imaged with the \emph{same} normal-FOV pinhole camera
($f_x = f_y = 1{,}649$~px, $c_x{=}960$, $c_y{=}540$, hFOV~${\approx}60^\circ$,
$1920\times1080$, zero distortion ground truth) to ensure a fair comparison.
Two calibration targets are studied, matched in corner count and physical extent
to enable fair cross-target comparisons:
\begin{itemize}
    \item \textbf{AprilGrid}: $6{\times}6$ tag grid (144 outer-corner points per
    view, $12{\times}12$ layout, 4.22\,cm pitch, $46.4{\times}46.4$\,cm corner span).
    \item \textbf{Matched ChArUco}: $13{\times}13$ square board (144 inner-corner points per
    view, $12{\times}12$ layout, 4.22\,cm pitch, $46.4{\times}46.4$\,cm corner span).
\end{itemize}
Both targets have identical corner count (144), pitch (4.22\,cm), and
$46.4\times46.4$\,cm corner span, enabling unconfounded cross-target comparison.
Both boards are generated analytically using the same camera parameters and
hemisphere geometry; noise is added as i.i.d.\ Gaussian per corner per view.
Three noise levels are evaluated: near-noiseless ($\sigma = 0.01$~px),
realistic ($\sigma = 0.5$~px), and high noise ($\sigma = 1.0$~px).
Distortion robustness experiments additionally apply medium and high
Brown--Conrady distortion (Section~\ref{subsec:distortion}).
The downstream calibration backend is \texttt{cv2.calibrateCamera} (OpenCV
Zhang-style) with a neutral focal-length initialization
$f_0 = 0.7\max(w,h)$ to avoid sensitivity to a bad starting point.

% -------------------------------------------------------
\subsection{Main Results: Fixed-Budget Subset Comparison}
\label{subsec:fixed_backend}

Tables~\ref{tab:aprilgrid_noiseless}--\ref{tab:charuco_sigma10} report the
five-method head-to-head comparison on both targets under three noise levels
($\sigma \in \{0,\,0.5,\,1.0\}$~px for AprilGrid; $\sigma{=}0.01$~px used
as near-noiseless for ChArUco).
All methods share the same candidate pool ($N{=}720$) and downstream
calibration backend; only the selected $K{=}20$ views differ.

\begin{table}[ht]
\centering
\setlength{\tabcolsep}{4pt}
\caption{AprilGrid, noiseless ($\sigma{\approx}0$\,px). $N{=}720$, $K{=}20$.
Focal/PP errors in pixels, distortion error unitless.
Random: mean over 25 trials. Bold: best E-optimal method (FW-E or Greedy-E).}
\label{tab:aprilgrid_noiseless}
\begin{tabular}{lrrrrrr}
\toprule
Method & $\lambda_{\min}$ & $\epsilon_f$\,(px) & $\epsilon_{pp}$\,(px) & $\epsilon_d$ & Train-RMS & Held-out \\
\midrule
\textbf{FW-E}     & $\mathbf{2.18\!\times\!10^{3}}$ & $\mathbf{4.11\!\times\!10^{-3}}$ & $1.88\!\times\!10^{-2}$ & $\mathbf{2.01\!\times\!10^{-4}}$ & $\mathbf{1.40\!\times\!10^{-2}}$ & $2.31\!\times\!10^{-2}$ \\
\textbf{Greedy-E} & $2.09\!\times\!10^{3}$ & $5.07\!\times\!10^{-3}$ & $\mathbf{1.51\!\times\!10^{-2}}$ & $6.00\!\times\!10^{-4}$ & $1.41\!\times\!10^{-2}$ & $\mathbf{1.98\!\times\!10^{-2}}$ \\
Random (mean)     & $892$ & $1.38\!\times\!10^{-2}$ & $2.17\!\times\!10^{-2}$ & $7.16\!\times\!10^{-4}$ & $1.39\!\times\!10^{-2}$ & $2.51\!\times\!10^{-2}$ \\
Coverage          & $761$ & $5.67\!\times\!10^{-3}$ & $3.13\!\times\!10^{-2}$ & $5.71\!\times\!10^{-4}$ & $1.40\!\times\!10^{-2}$ & $3.42\!\times\!10^{-2}$ \\
Motion-div.       & $663$ & $1.10\!\times\!10^{-2}$ & $2.47\!\times\!10^{-2}$ & $1.70\!\times\!10^{-4}$ & $1.38\!\times\!10^{-2}$ & $2.80\!\times\!10^{-2}$ \\
\bottomrule
\end{tabular}
\end{table}

\begin{table}[ht]
\centering
\setlength{\tabcolsep}{4pt}
\caption{Matched ChArUco board ($13{\times}13$, 144\,pts, 46\,cm), near-noiseless ($\sigma{=}0.01$\,px). $N{=}720$, $K{=}20$.
Focal/PP errors in pixels, distortion error unitless. Bold: best E-optimal method (FW-E or Greedy-E).}
\label{tab:charuco_noiseless}
\begin{tabular}{lrrrrrr}
\toprule
Method & $\lambda_{\min}$ & $\epsilon_f$\,(px) & $\epsilon_{pp}$\,(px) & $\epsilon_d$ & Train-RMS & Held-out \\
\midrule
FW-E              & $1.94\!\times\!10^{3}$ & $\mathbf{5.88\!\times\!10^{-3}}$ & $2.16\!\times\!10^{-2}$ & $1.63\!\times\!10^{-4}$ & $\mathbf{1.39\!\times\!10^{-2}}$ & $2.52\!\times\!10^{-2}$ \\
\textbf{Greedy-E} & $\mathbf{2.04\!\times\!10^{3}}$ & $8.43\!\times\!10^{-3}$ & $\mathbf{1.82\!\times\!10^{-2}}$ & $\mathbf{2.51\!\times\!10^{-5}}$ & $1.41\!\times\!10^{-2}$ & $\mathbf{2.25\!\times\!10^{-2}}$ \\
Random (mean)     & $885$ & $8.81\!\times\!10^{-3}$ & $1.41\!\times\!10^{-2}$ & $5.26\!\times\!10^{-4}$ & $1.36\!\times\!10^{-2}$ & $1.80\!\times\!10^{-2}$ \\
Coverage          & $715$ & $2.17\!\times\!10^{-2}$ & $4.92\!\times\!10^{-2}$ & $7.13\!\times\!10^{-4}$ & $1.38\!\times\!10^{-2}$ & $5.25\!\times\!10^{-2}$ \\
Motion-div.       & $592$ & $2.20\!\times\!10^{-2}$ & $6.12\!\times\!10^{-2}$ & $6.43\!\times\!10^{-4}$ & $1.40\!\times\!10^{-2}$ & $6.28\!\times\!10^{-2}$ \\
\bottomrule
\end{tabular}
\end{table}

\begin{table}[ht]
\centering
\setlength{\tabcolsep}{4pt}
\caption{AprilGrid, realistic noise ($\sigma{=}0.5$~px). $N{=}720$, $K{=}20$.
Focal/PP errors in pixels, distortion error unitless. Bold: best E-optimal method (FW-E or Greedy-E).}
\label{tab:aprilgrid_sigma05}
\begin{tabular}{lrrrrrr}
\toprule
Method & $\lambda_{\min}$ & $\epsilon_f$\,(px) & $\epsilon_{pp}$\,(px) & $\epsilon_d$ & Train-RMS & Held-out \\
\midrule
\textbf{FW-E}     & $\mathbf{0.874}$ & $0.203$ & $0.648$ & $0.0117$ & $\mathbf{0.695}$ & $0.904$ \\
\textbf{Greedy-E} & $0.842$ & $\mathbf{0.150}$ & $\mathbf{0.530}$ & $\mathbf{3.43\!\times\!10^{-3}}$ & $0.696$ & $\mathbf{0.830}$ \\
Random (mean)     & $0.359$ & $0.691$ & $1.066$ & $0.0358$ & $0.695$ & $1.237$ \\
Coverage          & $0.307$ & $0.278$ & $1.586$ & $0.0299$ & $0.702$ & $1.733$ \\
Motion-div.       & $0.268$ & $0.548$ & $1.244$ & $8.43\!\times\!10^{-3}$ & $0.691$ & $1.408$ \\
\bottomrule
\end{tabular}
\end{table}

\begin{table}[ht]
\centering
\setlength{\tabcolsep}{4pt}
\caption{Matched ChArUco board ($13{\times}13$, 144\,pts, 46\,cm), realistic noise ($\sigma{=}0.5$~px). $N{=}720$, $K{=}20$.
Focal/PP errors in pixels, distortion error unitless. Bold: best E-optimal method (FW-E or Greedy-E).}
\label{tab:charuco_sigma05}
\begin{tabular}{lrrrrrr}
\toprule
Method & $\lambda_{\min}$ & $\epsilon_f$\,(px) & $\epsilon_{pp}$\,(px) & $\epsilon_d$ & Train-RMS & Held-out \\
\midrule
FW-E              & $0.778$ & $\mathbf{0.293}$ & $\mathbf{1.089}$ & $8.81\!\times\!10^{-3}$ & $\mathbf{0.695}$ & $\mathbf{1.271}$ \\
\textbf{Greedy-E} & $\mathbf{0.809}$ & $0.424$ & $1.316$ & $\mathbf{5.01\!\times\!10^{-3}}$ & $0.700$ & $1.529$ \\
Random (mean)     & $0.356$ & $0.434$ & $0.698$ & $0.0259$ & $0.681$ & $0.898$ \\
Coverage          & $0.288$ & $1.077$ & $2.458$ & $0.0358$ & $0.692$ & $2.628$ \\
Motion-div.       & $0.239$ & $1.112$ & $3.104$ & $0.0328$ & $0.698$ & $3.183$ \\
\bottomrule
\end{tabular}
\end{table}

\begin{table}[ht]
\centering
\setlength{\tabcolsep}{4pt}
\caption{AprilGrid, $\sigma{=}1.0$~px noise. $N{=}720$, $K{=}20$.
Focal/PP errors in pixels, distortion error unitless. Bold: best E-optimal method (FW-E or Greedy-E).}
\label{tab:aprilgrid_sigma10}
\begin{tabular}{lrrrrrr}
\toprule
Method & $\lambda_{\min}$ & $\epsilon_f$\,(px) & $\epsilon_{pp}$\,(px) & $\epsilon_d$ & Train-RMS & Held-out \\
\midrule
\textbf{FW-E}     & $\mathbf{0.220}$ & $0.559$ & $1.994$ & $\mathbf{0.0364}$ & $\mathbf{1.398}$ & $2.404$ \\
\textbf{Greedy-E} & $0.217$ & $\mathbf{0.260}$ & $\mathbf{1.476}$ & $0.0526$ & $1.409$ & $\mathbf{2.002}$ \\
Random (mean)     & $0.0914$ & $1.377$ & $2.105$ & $0.0712$ & $1.390$ & $2.450$ \\
Coverage          & $0.0784$ & $0.551$ & $3.232$ & $0.0621$ & $1.405$ & $3.528$ \\
Motion-div.       & $0.0683$ & $1.093$ & $2.491$ & $0.0167$ & $1.383$ & $2.819$ \\
\bottomrule
\end{tabular}
\end{table}

\begin{table}[ht]
\centering
\setlength{\tabcolsep}{4pt}
\caption{Matched ChArUco board ($13{\times}13$, 144\,pts, 46\,cm), $\sigma{=}1.0$~px noise. $N{=}720$, $K{=}20$.
Focal/PP errors in pixels, distortion error unitless. Bold: best E-optimal method (FW-E or Greedy-E); FW-E collapses at this setting.}
\label{tab:charuco_sigma10}
\begin{tabular}{lrrrrrr}
\toprule
Method & $\lambda_{\min}$ & $\epsilon_f$\,(px) & $\epsilon_{pp}$\,(px) & $\epsilon_d$ & Train-RMS & Held-out \\
\midrule
FW-E              & $1.84\!\times\!10^{-2}$ & $8.34$ & $10.79$ & $0.0471$ & $\mathbf{1.382}$ & $10.52$ \\
\textbf{Greedy-E} & $\mathbf{0.204}$ & $\mathbf{0.944}$ & $\mathbf{3.243}$ & $\mathbf{0.0327}$ & $\mathbf{1.382}$ & $\mathbf{3.668}$ \\
Random (mean)     & $0.0907$ & $0.865$ & $1.392$ & $0.0516$ & $1.363$ & $1.792$ \\
Coverage          & $0.0738$ & $2.141$ & $4.913$ & $0.0716$ & $1.385$ & $5.256$ \\
Motion-div.       & $0.0614$ & $2.246$ & $6.261$ & $0.0673$ & $1.395$ & $6.417$ \\
\bottomrule
\end{tabular}
\end{table}

\paragraph{FIM methods substantially outperform geometric heuristics on AprilGrid.}
At $\sigma{=}0.5$\,px on AprilGrid (Table~\ref{tab:aprilgrid_sigma05}),
Greedy-E achieves the best $\epsilon_f{=}0.150$\,px, $\epsilon_{pp}{=}0.530$\,px,
and held-out RMS $0.830$\,px---improvements of $4.6{\times}$, $2.0{\times}$,
and $1.5{\times}$ over Random ($0.691$\,px, $1.066$\,px, $1.237$\,px).
FW-E leads $\lambda_{\min}$ ($0.874$ vs.\ Greedy-E $0.842$), confirming that
the Frank-Wolfe relaxation correctly solves the E-optimal objective, though
calibration quality favours the greedy discrete solution.
Coverage and Motion-diversity consistently produce the worst calibration
quality: held-out $1.733$ and $1.408$\,px respectively, confirming that
translational or rotational diversity heuristics do not target the most
informative views for intrinsic calibration.
The same pattern holds at $\sigma{=}1.0$\,px (Table~\ref{tab:aprilgrid_sigma10}):
Greedy-E wins $\epsilon_f$ ($0.260$\,px) and held-out RMS ($2.002$\,px)
versus Random $2.450$\,px, Coverage $3.528$\,px, and Motion-diversity $2.819$\,px.

\paragraph{ChArUco: Greedy-E is the recommended E-optimal method; FW-E collapses at high noise.}
On matched ChArUco, Greedy-E consistently achieves the highest $\lambda_{\min}$
($2$--$11{\times}$ above Random) and is the best-performing E-optimal method
across all noise levels.
At $\sigma{=}0.5$\,px (Table~\ref{tab:charuco_sigma05}), FW-E achieves the
best $\epsilon_f{=}0.293$\,px and held-out $1.271$\,px among E-optimal methods,
while Greedy-E leads $\lambda_{\min}$ ($0.809$) and $\epsilon_d$.
At $\sigma{=}1.0$\,px (Table~\ref{tab:charuco_sigma10}), Greedy-E dominates
all calibration metrics among E-optimal methods:
$\epsilon_f{=}0.944$\,px, $\epsilon_{pp}{=}3.243$\,px, $\epsilon_d{=}0.0327$,
and held-out $3.668$\,px.
FW-E collapses at this setting due to the rounding failure described below
($\lambda_{\min}{=}1.84{\times}10^{-2}$, $\epsilon_f{=}8.34$\,px).
Note that Random's lower absolute calibration errors at high noise reflect
ChArUco's more uniform inner-corner coverage incidentally constraining the
principal-point axis; however, Random provides no information certificate,
no $\lambda_{\min}$ guarantee, and is not a reliable strategy.
Greedy-E is the clear recommendation on ChArUco.

\paragraph{FW-E vs.\ Greedy-E: relaxation gap grows with noise.}
At low noise, FW-E and Greedy-E perform comparably on AprilGrid, with FW-E
winning $\lambda_{\min}$ and Greedy-E winning most calibration metrics.
This decoupling arises from top-$K$ rounding: the Frank-Wolfe continuous
solution is near-optimal for the relaxed objective, but rounding the
fractional weights to the top-20 integer indices introduces a gap that grows
as the information landscape becomes flatter at higher noise.
Greedy-E, which operates directly on discrete subsets at each step, avoids
this gap and maintains consistent calibration accuracy across noise levels.

\paragraph{FW-E collapses at high noise on ChArUco.}
At $\sigma{=}1.0$\,px on matched ChArUco (Table~\ref{tab:charuco_sigma10}),
FW-E collapses catastrophically: $\lambda_{\min}$ drops to
$1.84\!\times\!10^{-2}$ (an $11{\times}$ gap below Greedy-E's $0.204$),
$\epsilon_f$ rises to $8.34$\,px, and held-out RMS reaches $10.52$\,px---far
exceeding every other method including Random ($1.792$\,px).
The mechanism is the same rounding failure: the high-noise information
landscape is nearly flat, the continuous relaxation distributes weight
uniformly across candidates, and top-$K$ rounding assembles a near-random
set with low information content.
The FW duality gap certificate still closes (the relaxed problem converged),
but it no longer certifies the quality of the rounded discrete solution.

\paragraph{Train-RMS is not a reliable quality indicator.}
Across all methods and both targets, train-RMS lies in a narrow band
regardless of the wide variation in calibration errors.
At $\sigma{=}0.5$\,px the band is $0.69$--$0.70$\,px for AprilGrid and
$0.68$--$0.70$\,px for ChArUco; at $\sigma{=}1.0$\,px it is $1.36$--$1.41$\,px
for both targets, even as FW-E on ChArUco collapses to $\epsilon_{pp}{=}10.79$\,px.
Reprojection error on the training set is therefore an insufficient diagnostic
for calibration quality; held-out error and $\lambda_{\min}$ are necessary.

\paragraph{Coverage and Motion-diversity are consistently the weakest baselines.}
On ChArUco $\sigma{=}1.0$ (Table~\ref{tab:charuco_sigma10}), Motion-diversity
yields held-out $6.417$\,px and Coverage $5.256$\,px---$1.7{\times}$ and
$1.4{\times}$ worse than Greedy-E ($3.668$\,px) and $4.3{\times}$/$3.5{\times}$
worse on $\lambda_{\min}$.
On AprilGrid $\sigma{=}0.5$ (Table~\ref{tab:aprilgrid_sigma05}), Coverage
held-out ($1.733$\,px) is $2.1{\times}$ worse than Greedy-E ($0.830$\,px).
Maximising translational or rotational diversity does not maximise observability:
views spread across SO(3) do not necessarily provide the joint image-plane
coverage needed to constrain the full intrinsic parameter vector.

% -------------------------------------------------------
\subsection{FW-E Solver: Convergence and Runtime}
\label{subsec:fw_convergence}

The Frank-Wolfe continuous relaxation of~\eqref{eq:e_opt} exhibits strongly
noise-dependent convergence (Table~\ref{tab:fw_convergence}).
At low-to-moderate noise ($\sigma{\le}0.5$~px), FW wall-clock time is
$1.9$--$2.4$~s.
At high noise ($\sigma{=}1.0$~px), the information landscape is nearly flat
and the duality gap closes early; rounding the near-uniform continuous solution
degrades $\lambda_{\min}$ by $11\times$ on ChArUco
($1.84\!\times\!10^{-2}$ vs.\ Greedy-E $0.204$), driving held-out RMS
to $10.52$~px.
Greedy-E takes $\approx\!3.8$--$4.1$~s ($O(NK)$ Schur complements),
while geometric heuristics run in milliseconds.

\begin{table}[ht]
\centering
\caption{FW-E convergence summary ($N{=}720$, $K{=}20$).
Matched ChArUco uses $\sigma{=}0.01$\,px for the near-noiseless row.
At $\sigma{=}1.0$\,px (both targets), the high-noise landscape is nearly flat
and the duality gap closes early; rounding degrades $\lambda_{\min}$
significantly on ChArUco ($1.84\!\times\!10^{-2}$ rounded vs.\ Greedy-E $0.204$)
and $\epsilon_f$ rises to $8.34$\,px.}
\label{tab:fw_convergence}
\setlength{\tabcolsep}{5pt}
\begin{tabular}{llrrr}
\toprule
Target & $\sigma$ & Time (s) & Rounded $\lambda_{\min}$ & Held-out RMS\\
\midrule
ChArUco (matched) & 0.01 & 1.90 & $1.94\!\times\!10^{3}$  & $2.52\!\times\!10^{-2}$ \\
ChArUco (matched) & 0.5  & 1.86 & $7.78\!\times\!10^{-1}$ & $1.271$ \\
ChArUco (matched) & 1.0  & 1.20 & $1.84\!\times\!10^{-2}$ & $10.52$ \\
AprilGrid         & 0    & 2.40 & $2.18\!\times\!10^{3}$  & $2.31\!\times\!10^{-2}$ \\
AprilGrid         & 0.5  & 1.88 & $8.74\!\times\!10^{-1}$ & $0.904$ \\
AprilGrid         & 1.0  & 1.62 & $2.20\!\times\!10^{-1}$ & $2.404$ \\
\bottomrule
\end{tabular}
\end{table}

A key implementation fix was replacing the Linear Minimization Oracle (LMO)
from an LP solver (\texttt{scipy.linprog}) with a closed-form top-$K$
selection.
With FIM weights $\sim10^{18}$, the gradient magnitudes are $\sim10^{18}$,
causing \texttt{linprog} to stall indefinitely.
The LMO for the budget constraint $\sum u_i \leq K$, $u\in[0,1]^N$ has a
closed-form solution --- pick the $K$ indices with the most negative gradient
--- reducing per-iteration LMO cost from $O(N^3)$ (LP) to $O(N\log K)$.

% -------------------------------------------------------
\subsection{Correlation Between E-Optimality and Empirical Accuracy}
\label{subsec:correlation}

On AprilGrid $\sigma{=}0.5$, the ranking by $\lambda_{\min}$ reliably predicts
the ranking by held-out RMS: FW-E and Greedy-E ($\lambda_{\min}{=}0.874$ and
$0.842$) both achieve held-out RMS well below Random, Coverage, and
Motion-diversity ($1.237$--$1.733$\,px, Table~\ref{tab:aprilgrid_sigma05}).
On matched ChArUco, the E-optimal criterion continues to identify the best
discrete calibration strategy: Greedy-E achieves the highest $\lambda_{\min}$
at every noise level and is the top-performing E-optimal method on all
calibration metrics.
The absolute numbers on ChArUco are larger than AprilGrid because ChArUco's
inner-corner geometry provides more uniform image-plane coverage; incidentally,
Random's diverse sampling can better constrain the principal-point axis at high
noise---but this is an artifact of geometry rather than a principled strategy,
and Random carries no information certificate.
At high noise on ChArUco, FW-E's rounding failure decouples the
relaxed objective from the discrete solution ($\lambda_{\min}$ collapses
$11{\times}$, held-out RMS rises to $10.52$\,px), while Greedy-E's
direct discrete operation remains robust.
These results confirm Greedy-E as the recommended method on ChArUco and
motivate developing more robust rounding strategies for FW-E as future work.

% -------------------------------------------------------
\subsection{Distortion Robustness}
\label{subsec:distortion}

To test whether method rankings hold under lens distortion, we generated four
additional datasets using the same hemisphere geometry ($N{=}720$, $K{=}20$)
but with Brown--Conrady distortion coefficients applied analytically:
medium distortion ($k_1{=}{-}0.15$, $k_2{=}0.05$, $p_1{=}0.001$, $p_2{=}0.001$)
and high distortion ($k_1{=}{-}0.35$, $k_2{=}0.15$, $p_1{=}0.002$, $p_2{=}0.002$),
each for both targets.
Table~\ref{tab:distortion_medium} reports $\lambda_{\min}$ for the three methods;
Table~\ref{tab:distortion_high} reports the high-distortion results.

\begin{table}[ht]
\centering
\small
\caption{Medium distortion ($k_1{=}{-}0.15$, $k_2{=}0.05$, $p_1{=}p_2{=}0.001$).
$\lambda_{\min}(H_{\rm cal})$, $N{=}720$, $K{=}20$.
Bold: best per column within each target.}
\label{tab:distortion_medium}
\setlength{\tabcolsep}{4pt}
\begin{tabular}{lrrrrrr}
\toprule
 & \multicolumn{3}{c}{AprilGrid} & \multicolumn{3}{c}{Matched ChArUco} \\
\cmidrule(lr){2-4}\cmidrule(lr){5-7}
Method & $\sigma{\approx}0$ & $\sigma{=}0.5$ & $\sigma{=}1.0$ & $\sigma{=}0.01$ & $\sigma{=}0.5$ & $\sigma{=}1.0$ \\
\midrule
FW-E     & $\mathbf{2.41\!\times\!10^{3}}$ & $\mathbf{0.965}$ & $\mathbf{0.241}$ & $\mathbf{2.23\!\times\!10^{3}}$ & $\mathbf{0.899}$ & $\mathbf{0.225}$ \\
Greedy-E & $\mathbf{2.41\!\times\!10^{3}}$ & $0.926$ & $0.240$ & $2.17\!\times\!10^{3}$ & $0.865$ & $0.217$ \\
Random   & $937$ & $0.377$ & $0.096$ & $947$ & $0.381$ & $0.097$ \\
\bottomrule
\end{tabular}
\end{table}

\begin{table}[ht]
\centering
\small
\caption{High distortion ($k_1{=}{-}0.35$, $k_2{=}0.15$, $p_1{=}p_2{=}0.002$).
$\lambda_{\min}(H_{\rm cal})$, $N{=}720$, $K{=}20$.
${}^*$FW-E rounding failure.}
\label{tab:distortion_high}
\setlength{\tabcolsep}{4pt}
\begin{tabular}{lrrrrrr}
\toprule
 & \multicolumn{3}{c}{AprilGrid} & \multicolumn{3}{c}{Matched ChArUco} \\
\cmidrule(lr){2-4}\cmidrule(lr){5-7}
Method & $\sigma{\approx}0$ & $\sigma{=}0.5$ & $\sigma{=}1.0$ & $\sigma{=}0.01$ & $\sigma{=}0.5$ & $\sigma{=}1.0$ \\
\midrule
FW-E     & $\mathbf{3.53\!\times\!10^{3}}$ & $0.348^*$ & $0.089^*$ & $\mathbf{3.34\!\times\!10^{3}}$ & $\mathbf{1.34}$ & $0.059^*$ \\
Greedy-E & $3.48\!\times\!10^{3}$ & $\mathbf{1.39}$ & $\mathbf{0.349}$ & $3.31\!\times\!10^{3}$ & $1.32$ & $\mathbf{0.332}$ \\
Random   & $1.35\!\times\!10^{3}$ & $0.543$ & $0.137$ & $1.42\!\times\!10^{3}$ & $0.568$ & $0.144$ \\
\bottomrule
\end{tabular}
\end{table}

\paragraph{Medium distortion leaves rankings intact.}
Under medium distortion, FW-E wins $\lambda_{\min}$ on both targets at all
noise levels, and the $2.5{\times}$ FW-E/Random gap in $\lambda_{\min}$ is
preserved (Table~\ref{tab:distortion_medium}).
The project\_points Jacobian correctly accounts for the Brown--Conrady model,
so the FIM landscape shifts uniformly and optimal view placement remains
stable.

\paragraph{High distortion triggers FW-E rounding failure at moderate noise.}
Under high distortion ($k_1{=}{-}0.35$), FW-E fails at $\sigma{=}0.5$\,px on
AprilGrid ($\lambda_{\min}{=}0.348$ vs.\ Greedy-E $1.39$, a $4.0{\times}$ gap)
and at $\sigma{=}1.0$\,px on both targets
(AprilGrid: $0.089$ vs.\ $0.349$; ChArUco: $0.059$ vs.\ $0.332$).
Strong barrel distortion concentrates feature detection near the optical axis,
flattening the FIM landscape and triggering early rounding collapse even at
moderate noise---an amplification of the same mechanism observed under zero
distortion.
Greedy-E is the most robust method: it wins $\lambda_{\min}$ at $\sigma{=}0.5$
and $\sigma{=}1.0$ on both targets under high distortion and is the recommended
selection strategy when significant lens distortion is expected.

% -------------------------------------------------------
\subsection{Ablation Studies}
\label{subsec:ablations}

\paragraph{FIM-based selection vs.\ geometric heuristics across noise regimes.}
The advantage of E-optimal methods over geometric baselines is consistent on
AprilGrid but target-dependent on ChArUco
(Tables~\ref{tab:aprilgrid_noiseless}--\ref{tab:charuco_sigma10}):
\begin{itemize}
    \item \textbf{Near-noiseless:} FW-E wins $\lambda_{\min}$ on AprilGrid
    ($2.18\!\times\!10^3$); Greedy-E wins $\lambda_{\min}$ on ChArUco
    ($2.04\!\times\!10^3$) and $\epsilon_f$ ($5.88\!\times\!10^{-3}$\,px).
    All methods are well-conditioned; calibration quality differences are
    within solver-convergence noise.
    \item \textbf{$\sigma{=}0.5$~px:} FW-E wins $\lambda_{\min}$ on AprilGrid
    ($0.874$); Greedy-E leads calibration quality on AprilGrid
    ($\epsilon_f{=}0.150$\,px, held-out $0.830$\,px).
    On ChArUco, Random unexpectedly beats both FIM methods on held-out RMS
    ($0.898$\,px vs.\ FW-E $1.271$, Greedy-E $1.529$) and $\epsilon_{pp}$
    ($0.698$\,px vs.\ Greedy-E $1.316$\,px).
    \item \textbf{$\sigma{=}1.0$~px:} Greedy-E wins $\epsilon_f$ ($0.260$\,px)
    and held-out RMS ($2.002$\,px) on AprilGrid.
    On ChArUco, Random wins all calibration metrics while FW-E collapses
    ($\lambda_{\min}$ drops $11{\times}$, held-out $10.52$\,px).
\end{itemize}

\paragraph{Target type.}
Both targets are matched in corner count, pitch, and physical extent
(Section~\ref{sec:experiments}), so cross-target comparisons isolate
target geometry (inner-corner vs.\ outer-corner) from board size effects.
AprilGrid consistently achieves lower held-out RMS than matched ChArUco with
the best FIM method at each noise level: near-noiseless
($1.98\!\times\!10^{-2}$\,px vs.\ $2.25\!\times\!10^{-2}$\,px),
$\sigma{=}0.5$\,px ($0.830$\,px vs.\ $1.271$\,px on FW-E),
and $\sigma{=}1.0$\,px ($2.002$\,px vs.\ $3.668$\,px on Greedy-E).
The outer-corner layout of AprilGrid tags is more precisely localizable in
the projection model, yielding a better-conditioned FIM that correctly
guides view selection.
On ChArUco, the inner-corner geometry provides more uniform image-plane corner
sampling; the E-optimal criterion does not fully exploit this structure,
allowing Random selection to match or beat FIM methods on held-out accuracy.

\paragraph{Subset budget $K$.}
Table~\ref{tab:k_sweep} reports $\lambda_{\min}(H_{\mathrm{cal}})$ for three
methods across $K\in\{5,10,15,20,30,50,100,200\}$ on three
configurations ($N{=}720$ throughout).

\begin{table}[h]
\centering
\small
\caption{$\lambda_{\min}(H_{\mathrm{cal}})$ vs.\ budget $K$.
Top: ChArUco, $\sigma{=}0.5$\,px. Middle: AprilGrid, $\sigma{=}0.5$\,px.
Bottom: AprilGrid, $\sigma{\approx}0$\,px (near-noiseless). $N{=}720$.
Bold: best per column.
${}^*$FW-E rounding failure (relaxation too diffuse at small $K$).}
\label{tab:k_sweep}
\setlength{\tabcolsep}{4pt}
\begin{tabular}{lrrrrrrrr}
\toprule
Method & $K{=}5$ & $K{=}10$ & $K{=}15$ & $K{=}20$ & $K{=}30$ & $K{=}50$ & $K{=}100$ & $K{=}200$ \\
\midrule
\multicolumn{9}{l}{\textit{ChArUco, $\sigma=0.5$\,px}} \\[2pt]
Random   & $0.043$ & $0.122$ & $0.226$ & $0.354$ & $0.589$ & $0.906$ & $2.01$ & $4.03$ \\
Greedy-E & $\mathbf{0.197}$ & $0.415$ & $\mathbf{0.613}$ & $\mathbf{0.809}$ & $\mathbf{1.18}$ & $1.88$ & $\mathbf{3.54}$ & $\mathbf{6.63}$ \\
FW-E     & $0.002^*$ & $\mathbf{0.429}$ & $0.002^*$ & $0.778$ & $\mathbf{1.18}$ & $\mathbf{1.89}$ & $3.53$ & $\mathbf{6.63}$ \\
\midrule
\multicolumn{9}{l}{\textit{AprilGrid, $\sigma=0.5$\,px}} \\[2pt]
Random   & $0.045$ & $0.123$ & $0.224$ & $0.346$ & $0.589$ & $0.884$ & $2.00$ & $4.14$ \\
Greedy-E & $\mathbf{0.171}$ & $\mathbf{0.418}$ & $\mathbf{0.629}$ & $0.842$ & $1.25$ & $2.01$ & $3.74$ & $6.87$ \\
FW-E     & $0.001^*$ & $0.073^*$ & $0.235$ & $\mathbf{0.874}$ & $\mathbf{1.27}$ & $\mathbf{2.02}$ & $\mathbf{3.76}$ & $\mathbf{6.88}$ \\
\midrule
\multicolumn{9}{l}{\textit{AprilGrid, $\sigma{\approx}0$\,px}} \\[2pt]
Random   & $107$ & $301$ & $554$ & $859$ & $1.47\!\times\!10^{3}$ & $2.20\!\times\!10^{3}$ & $5.00\!\times\!10^{3}$ & $1.04\!\times\!10^{4}$ \\
Greedy-E & $\mathbf{428}$ & $\mathbf{1.04\!\times\!10^{3}}$ & $1.56\!\times\!10^{3}$ & $2.09\!\times\!10^{3}$ & $3.11\!\times\!10^{3}$ & $5.01\!\times\!10^{3}$ & $9.35\!\times\!10^{3}$ & $\mathbf{1.72\!\times\!10^{4}}$ \\
FW-E     & $0.001^*$ & $176$ & $\mathbf{1.62\!\times\!10^{3}}$ & $\mathbf{2.18\!\times\!10^{3}}$ & $\mathbf{3.17\!\times\!10^{3}}$ & $\mathbf{5.06\!\times\!10^{3}}$ & $\mathbf{9.38\!\times\!10^{3}}$ & $\mathbf{1.72\!\times\!10^{4}}$ \\
\bottomrule
\end{tabular}
\end{table}

\noindent\textbf{Greedy-E is the most reliable method across budgets.}
Greedy-E achieves the highest $\lambda_{\min}$ at $K{\le}15$ on both noisy
configurations and at $K{=}5$--$10$ on noiseless AprilGrid, where FW-E's
rounding fails.
At $K{\ge}20$ on AprilGrid, FW-E recovers and slightly exceeds Greedy-E
($0.874$ vs.\ $0.842$ at $K{=}20$, within $3.7\%$), and at large budgets the
two methods converge to nearly identical $\lambda_{\min}$.
Random uniformly underperforms both FIM methods across all tested budgets on
all three configurations.

\noindent\textbf{FW-E rounding instability at small $K$.}
FW-E fails at small $K$ on all three configurations.
On ChArUco $\sigma{=}0.5$, $\lambda_{\min}$ collapses to $0.002$ at $K{=}5$,
$100{\times}$ below Greedy-E ($0.197$); the failure recurs at $K{=}15$.
On AprilGrid $\sigma{=}0.5$, the failure persists through $K{=}10$
(FW-E $0.073$ vs.\ Greedy-E $0.418$, a $5.7{\times}$ gap) before recovering
at $K{\ge}15$.
On near-noiseless AprilGrid, FW-E collapses at $K{=}5$ (${\approx}0.001$ vs.\
Greedy-E $428$) and at $K{=}10$ ($176$ vs.\ Greedy-E $1040$), but leads from
$K{\ge}15$ onward.
At small $K$ ($K/N{\approx}0.007$--$0.021$), the initial uniform weights keep
the continuous relaxation near-flat, and top-$K$ rounding assembles a poor
discrete set; at $K{\ge}15$--$20$ the relaxation concentrates and rounding
succeeds.

\noindent\textbf{Diminishing returns.}
The marginal gain in $\lambda_{\min}$ per additional pose decreases as $K$
grows.
On ChArUco $\sigma{=}0.5$, Greedy-E delivers $\lambda_{\min}{=}0.197$ at $K{=}5$
but only ${\approx}0.010$ per additional pose in the $K{=}100\!\to\!200$ range
---an $18{\times}$ drop in per-pose efficiency.
The practical ``knee'' lies around $K{=}20$--$30$: beyond 30 frames, each
additional view adds less than $2\%$ of the cumulative gain already achieved.

% -------------------------------------------------------
\subsection{Real-World Calibration Study}
\label{subsec:real_world}

We validate the proposed framework on real ChArUco captures collected using
a live-demo rig (\texttt{live\_demo\_charuco/}).
Starting from an over-complete capture session, each method selects $K{=}20$
frames, and \texttt{cv2.calibrateCamera} is run on the selection.
Because no ground-truth intrinsics are available, we report held-out
reprojection error and cross-trial repeatability.
Results to be included in the final paper.

\printbibliography

\end{document}
